\paragraph{So sánh các chiến lược One-vs-Rest, One-vs-One và Softmax trực tiếp}

\tablename~kết quả~\ref{tab:ovr-ovo-softmax-compare-test} cho thấy sự khác biệt
rõ ràng giữa ba chiến lược phân loại đa lớp gồm One-vs-Rest (OvR), One-vs-One
(OvO) và Softmax trực tiếp.

\begin{table}[htbp]
  \centering
  \caption{Các chỉ số đánh giá của 3 chiến lược One-vs-Rest, One-vs-One và
  softmax trực tiếp trên tập kiểm tra của bộ dữ liệu Forest CoverType}
  \label{tab:ovr-ovo-softmax-compare-test}
  \import{\tablespath}{ovr-ovo-softmax-compare-test}

\end{table}

Về độ chính xác tổng thể (Accuracy), OvO Sigmoid đạt 0.7162, cao nhất nhưng chỉ
nhỉnh hơn Softmax (0.7148) và OvR (0.7073) không đáng kể. Tuy nhiên, khi xét
các chỉ số cân bằng hơn như Precision, Recall và F1-score macro, Softmax thể
hiện hiệu quả ổn định và tốt nhất với Precision bằng 0.5901, Recall bằng 0.4306
và F1-score bằng 0.4319.

OvO Sigmoid cho thấy cải thiện nhẹ về Recall so với OvR (0.4249 so với 0.3959),
đồng thời F1-score cũng cao hơn (0.4185 so với 0.3864), phản ánh khả năng cân
bằng tốt hơn giữa các lớp so với OvR. OvR lại có hiệu suất thấp nhất ở cả
Recall và F1-score, cho thấy hạn chế trong việc xử lý mất cân bằng giữa các lớp
khi huấn luyện theo từng lớp độc lập.

Nhìn chung, Softmax trực tiếp là lựa chọn ổn định hơn khi xét toàn bộ các chỉ
số macro, nhờ khả năng tối ưu hóa đồng thời xác suất trên tất cả lớp, giúp giảm
xung đột giữa các bộ phân loại nhị phân như trong OvR và OvO. OvO có ưu thế nhẹ
về Accuracy nhưng không vượt trội đáng kể, trong khi OvR tỏ ra kém hiệu quả
nhất trong bối cảnh bài toán này.

\figurename~\ref{fig:ovr-ovo-softmax-prec-cm} độ chính xác trên từng lớp và ma
trận nhầm lẫn trên tập kiểm tra của 3 chiến lược. Ta thấy rằng độ chính xác của
Softmax khá ổn định. Trong khi đó, OvO gặp khó ở lớp 5 với độ chính xác
bằng 0. Điều này cũng đã được thể hiện khi độ chính xác theo trung bình macro
cuả OvO là thấp nhất (0.564). Còn OvR bị gặp khó ở lớp 4 khi độ chính xác tụt
thấp hơn hẳn so với OvO và Softmax. Như nhìn chung, OvR có độ chính xác khá ổn
định giữa các lớp. Nhìn chung, ma trận nhầm lẫn giữa 3 chiến lược khá tương
đồng. OvR thì hay nhầm lẫn giữa lớp 7 thành 1 hơn. OvR và Softmax thì ít nhầm
lẫn đặc biệt hơn.

\begin{figure}[htbp]
  \centering
  \includegraphics[width=0.9\textwidth]{ovr-ovo-softmax-precision-confusion-matrix}
  \caption{Biểu đồ độ chính xác từng lớp và ma trận nhầm lẫn của 3 chiến lược
  OvR, OvO và Softmax trực tiếp trên tập kiểm tra của bộ dữ liệu Forest
  CoverType}
  \label{fig:ovr-ovo-softmax-prec-cm}
\end{figure}

Kết quả thực nghiệm cho thấy Softmax trực tiếp là phương pháp ổn định nhất khi
xét theo các chỉ số macro, trong khi OvO Sigmoid đạt Accuracy cao nhất nhưng
không vượt trội đáng kể và thiếu tính ổn định trên từng lớp. Đáng chú ý, OvO
không cải thiện hiệu năng trên các lớp khó như kỳ vọng, thậm chí còn thất bại
hoàn toàn ở một số lớp, cho thấy hạn chế của cách tiếp cận phân rã theo từng
cặp trong bối cảnh dữ liệu có sự chồng lấn đa lớp. Trong khi đó, OvR Sigmoid
cho kết quả tương đối ổn định nhưng mang tính bảo thủ, dẫn đến Recall và
F1-score thấp hơn.

Xét về mặt lý thuyết, OvR và OvO lần lượt yêu cầu $K$ và $K(K−1)/2$ mô hình,
trong khi Softmax chỉ cần một mô hình duy nhất, mang lại lợi thế rõ rệt về chi
phí tính toán. Tuy OvO có khả năng học ranh giới phân lớp chi tiết theo từng
cặp, kết quả thực nghiệm cho thấy lợi thế này không được phát huy trong trường
hợp cụ thể, có thể do cấu trúc dữ liệu không phù hợp với giả định phân rã cặp.

Nhìn chung, Softmax là lựa chọn phù hợp nhất cho bài toán này nhờ sự cân bằng
giữa hiệu năng và độ phức tạp.
